%% 
%% Copyright 2019-2024 Elsevier Ltd
%% 
%% This file is part of the 'CAS Bundle'.
%% --------------------------------------
%% 
%% It may be distributed under the conditions of the LaTeX Project Public
%% License, either version 1.3c of this license or (at your option) any
%% later version.  The latest version of this license is in
%%    http://www.latex-project.org/lppl.txt
%% and version 1.3c or later is part of all distributions of LaTeX
%% version 1999/12/01 or later.
%% 
%% The list of all files belonging to the 'CAS Bundle' is
%% given in the file `manifest.txt'.
%% 
%% Template article for cas-sc documentclass for 
%% double column output.

\documentclass[a4paper,fleqn]{cas-sc}

% If the frontmatter runs over more than one page
% use the longmktitle option.

%\documentclass[a4paper,fleqn,longmktitle]{cas-sc}

\usepackage[numbers,sort&compress]{natbib}
\usepackage[inkscapelatex=false]{svg}


%\usepackage[authoryear]{natbib}
%\usepackage[authoryear,longnamesfirst]{natbib}

%%%Author macros
\def\tsc#1{\csdef{#1}{\textsc{\lowercase{#1}}\xspace}}
\tsc{WGM}
\tsc{QE}
%%%

% Uncomment and use as if needed
%\newtheorem{theorem}{Theorem}
%\newtheorem{lemma}[theorem]{Lemma}
%\newdefinition{rmk}{Remark}
%\newproof{pf}{Proof}
%\newproof{pot}{Proof of Theorem \ref{thm}}

\begin{document}
\let\WriteBookmarks\relax
\def\floatpagepagefraction{1}
\def\textpagefraction{.001}

% Short title
\shorttitle{Short title, for the in-page naming}    

% Short author
\shortauthors{H. Álvarez, J. Carbó, J.M. Molina}  

% Main title of the paper
\title [mode = title]{Draft paper base bach thesis}  

% Title footnote mark
% eg: \tnotemark[1]
\tnotemark[1] 

% Title footnote 1.
% eg: \tnotetext[1]{Title footnote text}
\tnotetext[1]{} 

% First author
%
% Options: Use if required
% eg: \author[1,3]{Author Name}[type=editor,
%       style=chinese,
%       auid=000,
%       bioid=1,
%       prefix=Sir,
%       orcid=0000-0000-0000-0000,
%       facebook=<facebook id>,
%       twitter=<twitter id>,
%       linkedin=<linkedin id>,
%       gplus=<gplus id>]

\author[1]{H. Álvarez}[type=author,
        style=spanish,
        auid=000,
        bioid=1,
        orcid=0000-0000-0000-0000]

% Corresponding author indication
\cormark[1]

% Footnote of the first author
\fnmark[1]

% Email id of the first author
\ead{hectoral@pa.uc3m.es}

% URL of the first author
\ead[url]{}

% Credit authorship
% eg: \credit{Conceptualization of this study, Methodology, Software}
\credit{Aqui va la contribucion mia}



% Address/affiliation
\affiliation[1]{organization={Computer Science Department, Universidad Carlos III de Madrid},
            addressline={Av. Gregorio Peces-Barba Martínez 22}, 
            city={Colmenarejo},
%          citysep={}, % Uncomment if no comma needed between city and postcode
            postcode={28270}, 
            state={Madrid},
            country={Spain}}



\author[1]{J. Carbó}[type=author,
        style=spanish,
        auid=000,
        bioid=1,
        orcid=0000-0000-0000-0000]

% Footnote of the second author
\fnmark[1]

% Email id of the second author
\ead{jcarbo@inf.uc3m.es}

% URL of the second author
\ead[url]{}

% Credit authorship
\credit{Aqui va la contribucion de Javier/José Manuel}

\author[1]{J.M. Molina}[type=author,
        style=spanish,
        auid=000,
        bioid=1,
        orcid=0000-0000-0000-0000]

% Footnote of the second author
\fnmark[1]

% Email id of the second author
\ead{molina@ia.uc3m.es}

% URL of the second author
\ead[url]{}

% Credit authorship
\credit{Aqui va la otra contribucion de Javier/José Manuel}

% Address/affiliation
% \affiliation[2]{organization={},
%             addressline={}, 
%             city={},
% %          citysep={}, % Uncomment if no comma needed between city and postcode
%             postcode={}, 
%             state={},
%             country={}}

% Corresponding author text
%\cortext[1]{Corresponding author}

% Footnote text
\fntext[1]{}

% For a title note without a number/mark
%\nonumnote{}

% Here goes the abstract
\begin{abstract}
Here goes the abstract%% Remove this line from your manuscript.

% Lo pongo aqui no se muy bien por que. Creo que el enfoque adecuado para el paper debe ser el de,
% especificar que en casos donde los datos se presenten únicamente en forma de nº de carga y nº de
% descarga, mientras que la solucion general suele ser el representar al pasajero como agente
% e inferir su deseo mediante una OD matrix, al tener estas restricciones de datos junto 
% a que el objetivo no es analizar el comportamiento de los pasajeros, sino teniendo en cuenta el 
% número de pasajeros, analizar cómo se comportan los autobuses, un enfoque que no use OD, que se 
% suele estimar por IPF (con busqueda de maxima entropia y que DEPENDE MUCHO de la matriz semilla
% inicial, normalmente obtenida de encuestas o datos de tarjeta), se propone un enfoque donde la 
% carga y descarga se infiere de forma independiente. Esto se suma a, debido a que no se busca
% estudiar el comportamiento particular de pasajeros concretos, se puede reducir la granularidad
% al cargarnos el agente concreto, al no requerir de procesos coginitivos en él que alteren el
% comportamiento.
ABSTRAC ABSTRACT ABSTRACT
\end{abstract}

% Use if graphical abstract is present
%\begin{graphicalabstract}
%\includegraphics{}
%\end{graphicalabstract}

% Research highlights
\begin{highlights}
\item 
\item 
\item 
\end{highlights}


% Keywords
% Each keyword is seperated by \sep
\begin{keywords}
 \sep \sep \sep
\end{keywords}

\maketitle
\section{Introduction and literature review}\label{sec:intro}

Analysing public transport supply from observed operation can be a very complicated task for
governmental entities.
The quantities that determine whether capacity is
enough, such as peak loads and the occurrence of full buses, are extreme
values, and estimating them reliably requires a large number of observations
taken under the same conditions \cite{furth2006using}. The binding constraint,
however, is not the volume of data itself but the length of the period over
which the service configuration remains unchanged. Each modification of the supply
thus restarts the sample.
For this problem, multiple entities rely on simulation techniques to analyse
supply without real deployment \cite{toledo2010mesoscopic}.

This work focuses on the capacity of buses to absorb demand studying a digital simulation
of a bus route in Madrid, Spain, using a multi-agent system.


The multi-agent paradigm is a well-known and widely implemented technique in
transport simulation, and its use has increased in recent years
\cite{huang2022overview,divassonj2025abm,bastarianto2023agent}.
Despite this, in public transport, the main area of study in which it applies is
mostly real-time control systems \cite{balbo2024extended}, whereas this work focuses
on demand level analysis.

Some of the works among the literature of MAS in transportation take advantage of the
cognitive contribution certain paradigms
such as BDI (Belief-Desire-Intention) offer, like
modeling drivers and evacuees who must take decisions during a
simulation \cite{padgham2014integrating,singh2016integrating}.
Such architectures are appropriate when the decision process of each individual
agent is the object of study; when the focus is the aggregate behaviour of the
system, simpler agents are enough \cite{adam2016bdi}.
The architectural decision of which types of agents to use for a certain MAS depends on the
purpose the simulation aims at and on the set of
data available.


In terms of granularity and nature of the simulation, different types
appear in literature: macroscopic, mesoscopic and microscopic \cite{burghout2004hybrid}.
Microscopic simulations refer to those which describe the individual exact behaviour
of entities, with a level of detail up to second by second.
On the other hand, macroscopic simulations do not represent that individual behaviour,
but rather use a global perspective to obtain aggregated level indicators and without
the need to model invidividual entities.
In the middle of those appear the mesoscopic. While it is used to represent aggregated
level indicators, it still relies on using individual entities as so, but without the
need of having them modelled at detail \cite{toledo2010mesoscopic}.

The mesoscopic approach seems to suit better for the purpose of this work, as
each bus is represented as an entity, but its movement among the route
comes from sampled data from fitted distributions and does not change by traffic dynamic
inside the simulation.
A microscopic level of detail is something this kind of problem does
not require \cite{toledo2010mesoscopic}.

In this work passengers are used to analyse bus capacity, and the origin and destination of
individual entities are not required.
No external surveys or other data from which to construct OD matrices are
available. The usual approach is to infer the matrices via IPF with null seeding
\cite{furth1992bus,mccord2010iterative}, while in this work
the boarding and alighting of passengers is done via fitted distributions from real data operations.
This decision comes from the goal of the work. OD matrices assigns
origin and destination to passengers, but here the problem requires only to know how many boards
and alightings occurs in a specific stop given a bus arrival.
Several bus simulators already operate from stop-level counts without an OD
matrix. BusMezzo \cite{toledo2010mesoscopic} is a mesoscopic simulator in
which passengers are not represented individually. Arrivals at stops follow a
Poisson process, and the number of alighting passengers is drawn from those on
board. BusSim \cite{kieu2020dealing} is an agent-based model with bus and stop
agents in which passengers represents aggregate quantities, built on
synthetic data for short-term prediction. Talusan et al.
\cite{talusan2025dispatch} constructs an event-based simulator with APC
(Automatic Passenger Counting) data to evaluate dispatching decisions. These works establish
that occupancy can be studied without passenger destinations. However,
first and third works are not multi-agent systems, the second is not calibrated
on real operation, and none of them models the same line with and without
individual passengers.

For the agents used in this paper, as the work studies no individual decision 
and the available data contain
no information about one, no cognitive reasoning is represented for either
the bus or the passenger. The agents are therefore reactive in the classification
of Wooldridge and Jennings \cite{wooldridge1995intelligent}: they act in
response to events and to the messages they receive. They are not purely
reactive, however, as they keep an internal state
\cite{wooldridge1999intelligent}.
The form they have of interaction is through FIPA interaction protocols \cite{poslad2007specifying}.
The communication between different entities is the core of the behaviour of the
system.

From the engineering side, this work provides a multi-agent model of line E01
of EMT Madrid, calibrated with the boardings and alightings recorded by the
operator and used to analyse occupancy and demand absorption.
From the artificial intelligence side, it addresses how fine the agents of
such a system need to be. Two solutions are built that share the same bus
model and similar interaction protocols between reactive agents, and differ
only in whether each passenger is an agent or passengers are counts held by
stops and buses.
Comparing them isolates what the passenger agent contributes, in results and
in computational cost, when the available data record only counts.

The remainder of the paper is organised as follows.
Section~\ref{sec:solutions} presents the two proposed solutions, defining the
agents and the interaction protocols between them.
Section~\ref{sec:data} describes the operational data and how the boarding and
alighting distributions are calibrated from them.
Section~\ref{sec:results} details the implementation in GAMA, the simulation
parameters for line E01 and the results obtained.
Finally, Section~\ref{sec:conclusions} draws the conclusions.

\section{Proposed solutions: two different approaches}\label{sec:solutions}

Transport simulations usually represent each passenger as an agent, because a
passenger is an entity that makes choices \cite{cats2016dynamic}. The data
available here record only how many passengers board and alight at each stop,
so there is no individual choice to represent. Two solutions are therefore
built. The first solution keeps the passenger as an agent, as far as the data allow,
but the decision to alight is not its own. As a bus approaches a stop, the environment
draws the number of alighting passengers from the
fitted distribution and marks that many passengers on board, which then
request to alight when the bus asks. The second solution removes the passenger
agent, and buses and stops keep only the number of passengers. Comparing the
two shows what the passenger agent contributes when its behaviour cannot be
derived from the data.


\subsection{Agents and protocols}

As introduced in Section~\ref{sec:intro}, the agents modelled in this work are
reactive in the classification of Wooldridge and Jennings
\cite{wooldridge1995intelligent}. Their behaviour is a fixed set of rules
activated by events and by the messages they receive, with no deliberation.
They are, more specifically, agents with state \cite{wooldridge1999intelligent},
since each one keeps information between events, such as the occupancy of a
bus or the queue of passengers waiting at a stop, and acts according to it.
Each agent controls its own state and decides its response to a request
according to it; a bus, for instance, accepts or refuses a boarding request
depending on its remaining capacity.
This section defines the agents of the two solutions and the interaction
protocols through which they communicate.

In the two proposed implementations, three agents are defined:

\begin{itemize}
  \item Bus: Representing a phisical bus inside the simulation. It has a list of the trips it is
  scheduled to run, and keeps track of which and how many passengers has on board.
  It can present three different modes, which relates to its
  \textit{scheduled starting point}, the \textit{communications when its at a stop}, and
  the \textit{transit between stops}.
  \item Stop: Representing each stop of the bus route. It keeps track of which and how many passengers
  are waiting to the bus in what its considered its \textit{passengers custody} mode.
  When a bus arrives, it changes to \textit{communicating with a bus}.
  \item Passenger: This agent only appears in the first implementation. It represents each passenger
  individually, and has two possible modes:\textit{waiting at a stop} and \textit{on a bus}. 
  It has a moment of creation timestamp, which represents at which exact time it appears at a stop.
\end{itemize}

The agents interact through FIPA protocols, specifically Request \cite{fipa2002request} and Query
\cite{fipa2002query} types. Table \ref{tab:protocols} lists them in the order
in which they run each time a bus serves a stop. Since the agents share a
common vocabulary and their responses follow fixed rules, only the message
paths that can occur in the simulation are implemented: a request is always
agreed to, except in the boarding request.

In the first solution, as a bus approaches a stop it sends an \emph{alighting
query} to each passenger on board, which answers whether it will alight
according to the mark set by the environment. The bus then sends an
\emph{alighting request} to each passenger that answered affirmatively; the
passenger confirms and is removed from the simulation, and the bus updates its
list of passengers, which frees capacity. Next, the bus sends an \emph{arrival
notification} to the stop, and the stop sends a \emph{boarding notice} to each
waiting passenger, containing the reference of the bus to address. Each
passenger then sends a \emph{boarding request} to the bus, including the time
it has waited. This is the only protocol with a refusal: the bus agrees while
it has capacity and refuses once it is full, in which case the passenger
remains at the stop. Finally, the bus sends a \emph{count update} to the stop
with the passengers that boarded and their waiting times, which the stop
records.

In the second solution there are no passenger agents and two protocols remain.
On arrival, the bus sends an \emph{arrival query} to the stop, which returns
the number of passengers waiting and the number that must alight, the latter
obtained from the environment. The bus applies both figures against its
capacity and closes the visit with a \emph{count update} carrying the number
of passengers that boarded.

Every exchange consists of three messages, except a refused boarding request,
which takes two. In the first solution the number of messages per stop visit
therefore grows with the passengers on board and waiting, whereas in the
second it is constant at six.

\begin{table}[width=\linewidth,cols=4,pos=t]
\caption{Interaction protocols of the two solutions, in order of execution
during a stop visit.}\label{tab:protocols}
\begin{tabular*}{\tblwidth}{@{}p{.20\linewidth}p{.22\linewidth}p{.09\linewidth}p{.40\linewidth}@{}}
\toprule
Protocol & Initiator $\rightarrow$ participant & Type & Purpose \\
\midrule
\multicolumn{4}{@{}l}{\emph{First solution}} \\
Alighting query & Bus $\rightarrow$ passenger on board & Query &
  Whether the passenger alights at the next stop \\
Alighting request & Bus $\rightarrow$ alighting passenger & Request &
  The passenger leaves the bus and is removed \\
Arrival notification & Bus $\rightarrow$ stop & Request &
  The stop starts notifying its waiting passengers \\
Boarding notice & Stop $\rightarrow$ waiting passenger & Request &
  Tells the passenger which bus to address \\
Boarding request & Passenger $\rightarrow$ bus & Request &
  The bus agrees, or refuses when full \\
Count update & Bus $\rightarrow$ stop & Request &
  Who boarded and their waiting times \\
\midrule
\multicolumn{4}{@{}l}{\emph{Second solution}} \\
Arrival query & Bus $\rightarrow$ stop & Query &
  Number waiting and number alighting \\
Count update & Bus $\rightarrow$ stop & Request &
  Number of passengers that boarded \\
\bottomrule
\end{tabular*}
\end{table}


The FIPA diagrams that formalize this protocols are shown in Figures \ref{fig:},
\ref{fig:}, 

\newlength{\figh}

\begin{figure}
  \centering
  \setlength{\figh}{0.62\linewidth}
  \begin{minipage}[t]{0.764\figh}
    \centering
    \includesvg[height=\figh]{figs/protocols/en/consulta_de_descenso_en}\par
    \smallskip {\small (a) Alighting query}
  \end{minipage}\hspace{0.04\linewidth}
  \begin{minipage}[t]{0.676\figh}
    \centering
    \includesvg[height=\figh]{figs/protocols/en/diagrama_solicitud_de_descenso_en}\par
    \smallskip {\small (b) Alighting request}
  \end{minipage}
  \caption{FIPA-Request interaction protocols from first solution (1/3)}
  \label{fig:requests_first}
\end{figure}

\begin{figure}
  \centering
  \setlength{\figh}{0.62\linewidth}
  \begin{minipage}[t]{0.764\figh}
    \centering
    \includesvg[height=\figh]{figs/protocols/en/notificacion_de_llegada_en}\par
    \smallskip {\small (a) Arrival notification}
  \end{minipage}\hspace{0.04\linewidth}
  \begin{minipage}[t]{0.676\figh}
    \centering
    \includesvg[height=\figh]{figs/protocols/en/aviso_de_embarque_en}\par
    \smallskip {\small (b) Boarding notice}
  \end{minipage}
  \caption{FIPA-Request interaction protocols from first solution (2/3)}
  \label{fig:requests_seconds}
\end{figure}

\begin{figure}
  \centering
  \setlength{\figh}{0.62\linewidth}
  \begin{minipage}[t]{0.764\figh}
    \centering
    \includesvg[height=\figh]{figs/protocols/en/solicitud_de_embarque_en}\par
    \smallskip {\small (a) Boarding request}
  \end{minipage}\hspace{0.04\linewidth}
  \begin{minipage}[t]{0.676\figh}
    \centering
    \includesvg[height=\figh]{figs/protocols/en/actualizacion_recuento_en}\par
    \smallskip {\small (b) Count update}
  \end{minipage}
  \caption{FIPA-Request interaction protocols from first solution (3/3)}
  \label{fig:requests_seconds}
\end{figure}

\begin{figure}
  \centering
  \setlength{\figh}{0.62\linewidth}
  \begin{minipage}[t]{0.764\figh}
    \centering
    \includesvg[height=\figh]{figs/protocols/en/consulta_llegada_en}\par
    \smallskip {\small (a) Arrival query}
  \end{minipage}\hspace{0.04\linewidth}
  \begin{minipage}[t]{0.676\figh}
    \centering
    \includesvg[height=\figh]{figs/protocols/en/actualizacion_recuento_en}\par
    \smallskip {\small (b) Count update (duplicated)}
  \end{minipage}
  \caption{FIPA-Request interaction protocols from second solution}
  \label{fig:requests_seconds}
\end{figure}


\subsection{Environment}\label{sec:environment}
 
Some elements of the model do not belong to the domain of any agent and are
handled by the environment. Those are the passage of time,
the movement of buses between stops, the generation of passengers, the
selection of the passengers that alight, and the time a bus remains at a stop.
They are the same in both solutions, except where stated.
 
\paragraph{Time and movement}

Time advances in discrete steps of length $\Delta t$, from the start of the
service day until the last bus has completed all its scheduled trips. The
movement of a bus between two stops ignores spatial position and is purely
temporal. This process consists of a travel time to be completed and a counter that
increases by $\Delta t$ at each step. This avoids spatial computation and
reduces the computational load.
 
\paragraph{Generation of agents}

Bus and stop agents are created before the simulation starts: one stop agent
for each stop of each direction of the route, and one bus agent for each
physical vehicle, which receives the list of its scheduled departures and the
stop where its service begins. A bus agent is removed once it has completed
its service.
 
Passengers are generated during the simulation. At the start of each hour $h$,
a value $\tau_h$ is drawn for every stop from the boarding distribution of that
stop and time slot. Since that distribution is fitted to the boardings
aggregated over one hour, $\tau_h$ is the number of passengers that appear at
the stop during that hour. They are created progressively. At each step an
increment is added to a counter kept for the stop, and whenever its integer
part reaches one or more, that number of passengers is created and subtracted
from the counter. The increment is not constant within the hour but follows a
line that pivots on the midpoint of the hour and reaches, at its end, the
increment of the following hour:
\begin{equation}\label{eq:generation}
g_h(t) = \delta_h + \frac{2\,(\delta_{h+1}-\delta_h)}{T}
         \left(t - t_h - \frac{T}{2}\right),
\qquad t \in [t_h,\, t_h + T),
\end{equation}
where $t_h$ is the start of hour $h$, $T$ its duration, $N$ the number of steps
it contains and $\delta_h = \tau_h / N$. The mean of $g_h$ over the hour is
$\delta_h$, so the number of passengers generated in the hour is still
$\tau_h$; only the instants at which they appear are redistributed.

Note that generating passengers during hour $h$ requires $\tau_{h+1}$, the
number of passengers of the following hour. Each value is therefore sampled
one hour in advance. At the start of the simulation both $\tau_5$ and $\tau_6$
are sampled, where 5 and 6 are its first two time slots, and at the start of
every later hour $h$ the value sampled an hour earlier becomes $\tau_h$ and a
new $\tau_{h+1}$ is drawn.
 
In the first solution each generated passenger is an agent linked to its stop.
In the second, generation increases an integer counter at the stop, and the
instant of generation of each passenger is stored in a list, so that waiting
times can still be obtained.
 
\paragraph{Selection of alighting passengers}

The selection of the passengers that alight must correspond to the fitted
distributions. A destination cannot be assigned to a passenger when it is
created, because the stop at which it alights would depend on which bus
carries it and on the moment the bus reaches each stop; fixing it in advance
would be arbitrary. Moreover, any such assignment would make the number of
alightings per arrival an emergent quantity, sensitive to delays, bunching and
refusals for lack of capacity, and it would deviate from the alighting
distributions, which are fitted per arrival.
 
The selection is therefore made by the environment at the moment of alighting,
the only point at which an empirical distribution supports it. In the first
solution it is triggered by the bus when it starts the alighting query: a
value is drawn from the alighting distribution of the next stop and time slot,
and its integer part $n$ gives the number of passengers to alight. The $n$
passengers on board that waited longest at their stop are marked. The
criterion is deterministic, and choosing it over any other is irrelevant to
the purpose of the model, which is to reproduce the observed alighting counts.
The mark is the interface between the global selection and the reactive
behaviour of the passenger, which still carries out the corresponding
protocol. At the last stop of each direction all passengers alight.
 
In the second solution the stop invokes the same mechanism when it receives
the arrival query, and transmits the resulting number to the bus. In neither
solution does the knowledge behind the selection belong to an agent, they
only transmit its result.
 
\paragraph{Time at a stop}
In simulation steps, the exchange of messages at a stop takes the same time
whatever the number of passengers involved. The time a bus remains at a stop
is therefore extended to reflect the real duration of boarding and alighting,
under two premises. First, boarding and alighting take place in parallel, one
through each door, so the duration is bounded by the slower of the two flows
and not by their sum. Second, each passenger takes an approximately constant
time $t_p$. The time added at a stop is
\begin{equation}\label{eq:dwell}
E = \max\bigl( \max(b,\, a)\, t_p,\; E_{\min} \bigr),
\end{equation}
where $b$ and $a$ are the numbers of boarding and alighting passengers and
$E_{\min}$ is the time taken by the message exchanges, which is the minimum.
As a limitation, this minimum also applies at stops with no activity, where a
real bus would not stop; the bias is a small constant overhead per stop.
 
\subsection{Differences between the two solutions}\label{sec:differences}
 
Both solutions share the bus and stop agents, the schedule and capacity of the
buses, the movement between stops, the generation of passengers, the selection
of alighting passengers, the time at a stop and the recorded outputs.
Table~\ref{tbl:differences} summarises where they differ.
 
\begin{table}[width=\linewidth,cols=3,pos=t]
\caption{Differences between the two solutions.}\label{tbl:differences}
\begin{tabular*}{\tblwidth}{@{}>{\raggedright\arraybackslash}p{.24\linewidth}>{\raggedright\arraybackslash}p{.34\linewidth}>{\raggedright\arraybackslash}p{.34\linewidth}@{}}
\toprule
 & First solution & Second solution \\
\midrule
Passenger & Agent, from its creation at a stop until it alights &
  Not an agent; counted at stops and buses \\
Occupancy of a bus & List of the passenger agents on board &
  Number of passengers on board \\
State of a stop & List of the passenger agents waiting &
  Number of passengers waiting and list of their instants of generation \\
Role of a stop & Mediates between bus and passengers &
  Reports to the bus the passengers waiting and alighting \\
Alighting & Each marked passenger answers the bus &
  The stop transmits the number \\
Waiting time & Sent by each passenger in its boarding request &
  Obtained by the stop from the instants of generation \\
Protocols & Six & Two \\
Messages per stop visit & Grow with the passengers involved & Six \\
\bottomrule
\end{tabular*}
\end{table}
 
\subsection{Recorded outputs}\label{sec:outputs}
 
Both solutions record the same data, so that the indicators are obtained in
the same way. Each time a bus serves a stop, the following are stored: the
passengers that boarded, the passengers that could not board because the bus
was full, the passengers that alighted, the occupancy of the bus after the
operation, and the travel time from the previous stop. The waiting time of
each passenger that boarded is stored as well. Departures from the first stop
that take place later than scheduled are recorded as delays.

The complete outputs are recomputed into defined Key Performance Indicators (KPIs)
for demand absorption measurement.

This indicators can be divided into three main categories, depending of the scope
they are aimed at: general values, route values, and trip values.

The general values presents the following indicators with its comptue method






\section{Data calibration}\label{sec:data}

\subsection{Operational data}\label{sec:data_source}

The simulator is built on the operational records of line E01 provided by
EMT Madrid, covering an eight-month window starting July 2025.
The records follow a two-table model linked by a trip identifier. Each
departure from a terminal is an entry of the first table, with its date,
time, time slot, direction, vehicle and type of day. Each arrival of that
trip at a stop is an entry of the second table, with the time of arrival, its
time slot, and the number of passengers that boarded and alighted. The
occupancy after each stop is computed as the previous occupancy plus the
boardings minus the alightings.

Cleaning focus on the boarding, alighting and occupancy fields, where
negative values and extreme values derived from APC recording errors were found.
The cleaning also involved the removal of any recorded occupancy value above
a stablished threshold of 80 passengers.
The removed records amount, on average, to one record
every two days of data at each stop.

The study is restricted to normal days, defined as working days without free
fare. Weekends and public holidays are excluded because the window starts in
July 2025, so their number of observations is small compared with working
days, and parameters fitted on them would be biased by the small sample. Days
with free fare are excluded because the outlier analysis showed that their
occupancy is systematically higher than on the remaining working days, and
they add noise to the distributions. The window does not cover a full annual
cycle, so seasonal effects cannot be separated from the data, and the fitted
parameters describe the observed period.

\subsection{Calibration of the inputs}\label{sec:data_inputs}

The simulator receives four inputs: the boardings and alightings at each stop,
the travel times between consecutive stops, and the schedule of the buses.
The first three are fitted per time slot, using the slot of the moment at
which the bus serves the stop; the schedule is taken directly from the
operator, as the list of departures of each physical bus.

Boardings and alightings are modelled separately, and no origin-destination
matrix is estimated (Section~\ref{sec:intro}). They are also modelled at
different levels, according to how they are used. Boardings are aggregated
per stop and hour, since the simulator uses them as the number of passengers
that appear at the stop during that hour
(Section~\ref{sec:environment}). Alightings are modelled per bus arrival,
since that is the moment at which the simulator draws them.

Unlike occupancy, boardings and alightings are not normal in all cases: many
stop and slot combinations show values close to zero and long right tails.
Three candidate distributions are therefore considered: normal, negative
binomial and geometric. For each stop and time slot, outliers are first
removed with the median absolute deviation, which does not assume any
particular distribution \cite{leys2013detecting}; a goodness-of-fit test is
then run for each candidate, and the candidate with the lowest Akaike
information criterion \cite{akaike1974new} among those not rejected is
selected. When all three are rejected, the one with the lowest criterion is
kept. Overall, the normal distribution is selected in 62.4\% of the cases, the
negative binomial in 23.9\%, more often for boardings than for alightings, and
the geometric in the remaining 13.7\%.

Travel times between consecutive stops showed signs of normality from the
first visual inspection, so they are fitted only to normal distributions, per
segment and time slot. Table~\ref{tbl:nonrejection} gives the share of stop
or segment and time slot combinations whose fitted distribution is not
rejected.

\begin{table}[width=.8\linewidth,cols=3,pos=t]
\caption{Share of combinations whose fitted distribution is not rejected, per
direction of the line.}\label{tbl:nonrejection}
\begin{tabular*}{\tblwidth}{@{}LLL@{}}
\toprule
Variable & Towards Cibeles & Towards La Peseta \\
\midrule
Boardings & 90.60\% & 87.56\% \\
Alightings & 93.61\% & 90.91\% \\
Travel times & 95.49\% & 90.43\% \\
Occupancy (validation) & 94.39\% & 88.60\% \\
\bottomrule
\end{tabular*}
\end{table}

Since alightings are drawn per bus arrival and independently of the number of
passengers on board, the number of alightings over a period grows with the
number of buses that serve the stop. The parameters are therefore tied to the
service observed in the data, and the model reproduces that service; it is
not suited to evaluating changes of frequency without recalibration.

\subsection{Reference for validation}\label{sec:data_validation}

Occupancy is not an input of the simulator. It results from the combination
of boardings, alightings and the number of buses. It is used after the
simulation as the reference to validate the model. Its normality was studied
per stop and time slot with a test based on skewness and kurtosis at a
significance level of 0.05 \cite{dagostino1973tests,kim2013normality}, with outliers
detected by the generalised ESD test \cite{rosner1983}. Treating all
working days together gives a higher share of non-rejected combinations than
treating each weekday separately, so no distinction by weekday is made in the
rest of the work. For this variable, the time slot is that of the departure of
the trip from its terminal. Rejections concentrate in the last time slot of
the day, which has the fewest buses and therefore the fewest observations.

The terminal stops are excluded from this reference. It was reported that the
transition between consecutive trips is not always recorded at the right
moment, so the passengers boarding the next trip can be counted as occupancy
at the last stop of the previous one, and vice versa.

\subsection{Parametric implementation}\label{sec:data_sampling}

The fitted parameters are loaded before the simulation starts and updated
every hour of simulated time. Values are drawn with different methods for each family
according to what the platform offers. Normal values use the native generator
of GAMA, based on the Box--Muller transform \cite{box1958}.
The negative binomial is sampled as a Gamma--Poisson mixture
\cite{devroye1986}, which is valid for non-integer values of $r$:
\begin{equation}\label{eq:negbin}
\lambda \sim \mathrm{Gamma}\!\left(r,\ \frac{1-p}{p}\right), \qquad
X \mid \lambda \sim \mathrm{Poisson}(\lambda),
\end{equation}
with the Gamma distribution in shape and scale form. The geometric
distribution has a closed-form inverse:
\begin{equation}\label{eq:geom}
X = \left\lfloor \frac{\ln U}{\ln(1-p)} \right\rfloor, \qquad
U \sim \mathcal{U}(0,1),
\end{equation}
with support $\{0, 1, 2, \dots\}$, consistent with the parametrisation used in
the fitting.


\section{Implementations and results}\label{sec:results}

Both solutions are implemented in GAMA \cite{taillandier2019gama}, a platform
dedicated to agent-based modelling and simulation. It was chosen for two of
its features. Its FIPA skill provides each agent with a mailbox and supports
the FIPA communicative acts and interaction protocols, as well as
conversations whose sequence is defined by the modeller \cite{gama_fipa}. Its
headless mode runs experiments from the command line without the graphical
interface, so the replicas can be executed in batch \cite{gama_headless}.


\subsection{Case study}

For the simulation characterization, the layout of route Express 01 is used. The route
operates between the central square \textit{Cibeles}, located in the city center of the Spanish capital,
and \textit{La Peseta}, located in the southern neibourhood of Carabanchel.
The route has 15 and 12 stops for each direction respectively, and its core behavour reflects
the day by day of madrid's workers: Going to the center earlier in the morning to work and coming back home after lunch.
The route's prime time in Cibeles's direction comes between 7 and 9 in 
the morning, boarding passengers in the first stops of the trip, and alighting them at the final ones.
The same behaviour is inversed in La Peseta's direction, presenting its prime time after 3 pm, 
capturing the tendency of workers to return to its homeplace.

\subsection{Simulation setup}

Table~\ref{tbl_sim_params} lists the parameters of the experiments. The time
step $\Delta t$, the time per passenger $t_p$ and the minimum time at a stop
$E_{\min}$ are those of Eq.~(\ref{eq:dwell}), and the capacity of the buses
is the one adopted in Section~\ref{sec:data_source}.
 
\begin{table}[width=.9\linewidth,cols=2,pos=h]
\caption{Simulation parameters.}\label{tbl_sim_params}
\begin{tabular*}{\tblwidth}{@{}LL@{}}
\toprule
Parameter & Value \\
\midrule
Line & E01 (Cibeles--La Peseta) \\
Physical buses & 12 \\
Scheduled trips per day & 163 \\
Bus capacity & 80 passengers \\
Start of the simulation & 05:00 \\
Time step ($\Delta t$) & 2 s \\
Time per passenger ($t_p$) & 5 s \\
Minimum time at a stop ($E_{\min}$) & $4\,\Delta t$ \\
Replicas per solution & 1000 \\
\bottomrule
\end{tabular*}
\end{table}

% TODO 1 (execution, 2-3 sentences): how the replicas were run: headless
% batch experiment, one random seed per replica, how many replicas at once,
% and that both solutions ran on the same machine under the same limits.
% Give the machine (CPU, cores, memory) and any container limits: the run
% times compared in the results mean nothing without them.
% TODO 2 (what is measured, 1 sentence): the outputs of
% Section~\ref{sec:outputs}, plus the wall-clock time of each replica and,
% if you add it, the number of messages exchanged per replica.


\subsection{Evaluation}\label{sec:evaluation}
 
% TODO 3 (validation, 2-3 sentences): how each solution is compared with the
% observed occupancy of Section~\ref{sec:data_validation}: per stop and time
% slot, which two distributions are compared, with which measure (the
% Bhattacharyya coefficient in the thesis), and the acceptance threshold.
% State what supports the threshold you use.
% TODO 4 (solution against solution, 1-2 sentences): the measure of the
% difference between the two solutions per stop and time slot, and what it
% is compared against.
% TODO 5 (cost, 1 sentence): run time per replica and, if counted, messages.
 




\section{Conclusions}\label{sec:conclusions}
% Numbered list
% Use the style of numbering in square brackets.
% If nothing is used, default style will be taken.
%\begin{enumerate}[a)]
%\item 
%\item 
%\item 
%\end{enumerate}  

% Unnumbered list
%\begin{itemize}
%\item 
%\item 
%\item 
%\end{itemize}  

% Description list
%\begin{description}
%\item[]
%\item[] 
%\item[] 
%\end{description}  

%\clearpage %%Remove this from your manuscript


% % Figure
% \begin{figure}%[]
%   \centering
% %    \includegraphics{}
%     \caption{}\label{fig1}
% \end{figure}


% \begin{table}%[]
% \caption{}\label{tbl1}
% \begin{tabular*}{\tblwidth}{@{}LL@{}}
% \toprule
%   &  \\ % Table header row
% \midrule
%  & \\
%  & \\
%  & \\
%  & \\
% \bottomrule
% \end{tabular*}
% \end{table}







% Uncomment and use as the case may be
%\begin{theorem} 
%\end{theorem}

% Uncomment and use as the case may be
%\begin{lemma} 
%\end{lemma}

%% The Appendices part is started with the command \appendix;
%% appendix sections are then done as normal sections
%% \appendix

%\section{}\label{}

% To print the credit authorship contribution details
\printcredits

%% Loading bibliography style file
\bibliographystyle{elsarticle-num-names}
%\bibliographystyle{cas-model2-names}

% Loading bibliography database
\bibliography{cas-refs}

% Biography
%\bio{}
% Here goes the biography details.
%\endbio

%\bio{pic1}
% Here goes the biography details.
%\endbio

\end{document}

